\section{Projeto Integrador em Inteligência Artificial}

\textbf{OBJETIVO GERAL:}

Desenvolver competências para analisar problemas reais, projetar, implementar, validar, implantar e documentar soluções computacionais utilizando Inteligência Artificial, Machine Learning, IA Generativa, Engenharia de Dados, IoT e boas práticas de Engenharia de Software, integrando os conhecimentos desenvolvidos ao longo do curso.

\textbf{CARGA HORÁRIA:} 180 horas

\textbf{FUNÇÃO:}

Desenvolver sistemas computacionais baseados em Inteligência Artificial, Análise Preditiva e Ecossistemas IoT, atendendo às normas de qualidade corporativas, robustez metodológica, governança de dados e arquitetura segura.

\textbf{SUBFUNÇÕES:}

\begin{itemize}
\item Planejar soluções computacionais baseadas em Inteligência Artificial.
\item Desenvolver aplicações completas utilizando tecnologias estudadas ao longo do curso.
\item Integrar modelos de IA, bancos de dados, APIs e aplicações.
\item Validar técnica e funcionalmente as soluções desenvolvidas.
\item Documentar e apresentar projetos tecnológicos.
\end{itemize}

\textbf{PADRÕES DE DESEMPENHO:}

\begin{itemize}
\item Levantar requisitos junto ao cliente ou cenário proposto.
\item Planejar a arquitetura da solução.
\item Desenvolver aplicação funcional.
\item Integrar modelos de Inteligência Artificial.
\item Aplicar práticas de MLOps durante o desenvolvimento.
\item Validar desempenho da solução.
\item Elaborar documentação técnica completa.
\item Apresentar o projeto demonstrando domínio técnico.
\end{itemize}

\textbf{CAPACIDADES TÉCNICAS:}

\textbf{Planejamento}

\begin{itemize}
\item Levantar requisitos do projeto.
\item Definir objetivos técnicos.
\item Elaborar cronograma de desenvolvimento.
\item Planejar arquitetura da solução.
\end{itemize}

\textbf{Desenvolvimento}

\begin{itemize}
\item Desenvolver aplicações Full Stack.
\item Integrar APIs e modelos de IA.
\item Construir pipelines de dados.
\item Implementar funcionalidades inteligentes.
\end{itemize}

\textbf{Integração}

\begin{itemize}
\item Integrar Machine Learning às aplicações.
\item Integrar modelos generativos.
\item Integrar bancos relacionais, NoSQL e vetoriais.
\item Integrar serviços Cloud.
\end{itemize}

\textbf{Validação}

\begin{itemize}
\item Executar testes funcionais.
\item Avaliar desempenho da solução.
\item Corrigir falhas identificadas.
\item Validar requisitos do projeto.
\end{itemize}

\textbf{Entrega}

\begin{itemize}
\item Elaborar documentação técnica.
\item Produzir documentação do usuário.
\item Apresentar solução desenvolvida.
\item Defender tecnicamente o projeto.
\end{itemize}

\textbf{CONHECIMENTOS:}

\textbf{Gestão de Projetos}

\begin{itemize}
\item Planejamento
\item Cronograma
\item Escopo
\item Gestão de Riscos
\item Gestão de Requisitos
\end{itemize}

\textbf{Arquitetura de Software}

\begin{itemize}
\item Arquitetura em Camadas
\item Microsserviços
\item APIs REST
\item Integração de Sistemas
\item Arquitetura Baseada em Eventos
\end{itemize}

\textbf{Integração de Inteligência Artificial}

\begin{itemize}
\item Machine Learning
\item Deep Learning
\item IA Generativa
\item Grandes Modelos de Linguagem
\item Agentes Inteligentes
\item RAG
\end{itemize}

\textbf{Engenharia de Dados}

\begin{itemize}
\item ETL
\item ELT
\item PostgreSQL
\item MongoDB
\item Bancos Vetoriais
\end{itemize}

\textbf{Infraestrutura}

\begin{itemize}
\item Docker
\item Kubernetes
\item Cloud Computing
\item Deploy
\item Versionamento
\end{itemize}

\textbf{Qualidade}

\begin{itemize}
\item Testes Unitários
\item Testes de Integração
\item Testes Funcionais
\item Validação de Modelos
\item Documentação
\end{itemize}

\textbf{MLOps}

\begin{itemize}
\item MLflow
\item Versionamento
\item Monitoramento
\item Observabilidade
\item Governança
\end{itemize}

\textbf{Apresentação Técnica}

\begin{itemize}
\item Comunicação Técnica
\item Relatórios
\item Demonstração do Sistema
\item Pitch Técnico
\item Defesa do Projeto
\end{itemize}

\textbf{CAPACIDADES SOCIOEMOCIONAIS:}

\begin{itemize}
\item Demonstrar autonomia na resolução de problemas complexos.
\item Trabalhar colaborativamente em equipes multidisciplinares.
\item Liderar atividades técnicas quando necessário.
\item Comunicar soluções técnicas de forma clara e objetiva.
\item Produzir documentação técnica consistente.
\item Demonstrar pensamento crítico na tomada de decisões.
\item Gerenciar tempo e prioridades durante o desenvolvimento do projeto.
\item Agir com ética, responsabilidade e comprometimento profissional.
\item Adaptar-se a mudanças de requisitos e tecnologias.
\item Buscar melhoria contínua da solução desenvolvida.
\end{itemize}

\textbf{AMBIENTES PEDAGÓGICOS:}

\begin{itemize}
\item Laboratório de Inteligência Artificial.
\item Laboratório de Desenvolvimento de Software.
\item Laboratório de Redes e IoT.
\item Ambiente Virtual de Aprendizagem.
\item Ambiente Cloud para desenvolvimento e implantação.
\end{itemize}

\textbf{EQUIPAMENTOS E RECURSOS:}

\textbf{Hardware}

\begin{itemize}
\item Computadores com acesso à Internet.
\item Servidores para implantação das aplicações.
\item Cluster para execução dos modelos de IA.
\item Projetor multimídia.
\end{itemize}

\textbf{Software}

\begin{itemize}
\item Python
\item VS Code
\item Git
\item Docker
\item Kubernetes
\item PostgreSQL
\item MongoDB
\item ChromaDB
\item FastAPI
\item React
\item LangGraph
\item CrewAI
\item MLflow
\item Grafana
\item Jupyter Notebook
\end{itemize}
